\section{An exact central-section reduction by finite trapezoids}
The four centrally symmetric Platonic solids admit a common reduction:
for each plane orientation, a central section has at least as much area
as every parallel section.  In the present three-dimensional problem
this needs only a finite polygon construction.  A general
Brunn--Minkowski theorem, a limiting argument, and integration are
unnecessary.

\subsection{The planar finite construction}
Let $S$ be a convex polygon with at most $m$ vertices.  We prove
\[
 \operatorname{area}\bigl((S-S)/2\bigr)
 \ge \operatorname{area}(S).
\]
An empty polygon or a polygon of area zero causes no difficulty: its
area is zero, and the desired inequality follows from nonnegativity.
For completeness the construction below also handles a nonvertical
segment.  Let the horizontal projection be $[a,b]$.  If $a=b$, the
polygon lies on a vertical line and has area zero.  Assume $a<b$.

Write its vertical fibers as $[\ell(x),u(x)]$ for $a\le x\le b$.
The functions $\ell$ and $u$ are respectively convex and concave,
continuous piecewise-affine functions.  Their breakpoints occur among
the horizontal coordinates of the original polygon vertices.  This
follows directly by reading the lower and upper boundary chains of a
convex polygon.  Put $f(x)=u(x)-\ell(x)\ge0$, and let
\[
 c=(a+b)/2,\qquad \rho(x)=a+b-x.
\]
Construct a new polygon $W$ whose fiber over horizontal coordinate
$x-c$ has endpoints
\[
 L(x)=\frac{\ell(x)-u(\rho(x))}{2},\qquad
 H(x)=\frac{u(x)-\ell(\rho(x))}{2}.
\]
Their difference is
\[
 h(x)=H(x)-L(x)
      =\frac{f(x)+f(\rho(x))}{2}\ge0.                 \tag{CS1}
\]
The lower function $L$ is convex, and the upper function $H$ is
concave: composition with the affine reflection $\rho$ preserves
convexity or concavity, and taking the negative interchanges them.
Thus the region between these boundary chains is a convex polygon.

Its containment is also explicit.  The point $(x-c,L(x))$ is the
half-difference of $(x,\ell(x))\in S$ and
$(\rho(x),u(\rho(x)))\in S$, while $(x-c,H(x))$ is the
half-difference of $(x,u(x))\in S$ and
$(\rho(x),\ell(\rho(x)))\in S$.  Indeed their first coordinates
are $(x-\rho(x))/2=x-c$.  Since $(S-S)/2$ is convex, it contains the
whole segment between these endpoints and consequently all of $W$.
Equivalently, it suffices to check the finitely many boundary vertices
of $W$ and then take their convex hull.

\subsection{Exact equality of areas}
Take the distinct horizontal vertex coordinates of $S$, adjoin their
reflections under $\rho$, and sort the resulting finite set:
\[
 a=x_0<x_1<\cdots<x_r=b.
\]
There are at most $2m$ distinct breakpoints, so $r\le2m-1$.
The list is reflection invariant:
\[
 x_i+x_{r-i}=a+b \qquad(0\le i\le r).
\]
Between consecutive breakpoints all four functions
$\ell,u,\ell\circ\rho,u\circ\rho$ are affine.  Hence $S$ is the
union, with disjoint interiors, of its vertical trapezoids over these
intervals.  If $f_i=f(x_i)$ and $\Delta_i=x_{i+1}-x_i$, their areas are
\[
 T_i=\frac{\Delta_i}{2}(f_i+f_{i+1}),\qquad
 \operatorname{area}(S)=\sum_{i=0}^{r-1}T_i.          \tag{CS2}
\]
A zero endpoint height merely gives a triangle.  The trapezoid formula
follows by splitting into triangles or by a four-vertex shoelace
calculation; finite area additivity gives the second equality.

The corresponding slab of $W$ runs from $x_i-c$ to $x_{i+1}-c$.
Its horizontal width is still exactly $\Delta_i$: the coordinate map
$x\mapsto x-c$ is a translation, not a dilation.  From (CS1), its
endpoint heights are $(f_i+f_{r-i})/2$ and
$(f_{i+1}+f_{r-i-1})/2$.  Its area is therefore
\[
 \begin{split}
 U_i
 &=\frac{\Delta_i}{4}
      (f_i+f_{i+1}+f_{r-i}+f_{r-i-1})\\
 &=\frac12 T_i+\frac12 T_{r-i-1},
 \end{split}                                       \tag{CS3}
\]
because reflection gives $\Delta_i=\Delta_{r-i-1}$.
Summing (CS3), the reversed finite index list is a permutation of the
original list, and we obtain
\[
 \operatorname{area}(W)
 =\sum_{i=0}^{r-1}U_i
 =\frac12\sum_{i=0}^{r-1}T_i
  +\frac12\sum_{i=0}^{r-1}T_{r-i-1}
 =\operatorname{area}(S).                           \tag{CS4}
\]
As $W\subseteq(S-S)/2$, area monotonicity proves the asserted inequality.
For polygons, this monotonicity is itself a finite fact: refine by the
finitely many vertical breakpoints of both polygons, compare the
nonnegative fiber heights at the endpoints, and use the same trapezoid
formula.  No measure-theoretic premise is needed.

\subsection{Application to an affine section of a symmetric polytope}
Suppose $K=-K$ is a convex three-dimensional polytope, and write
\[
 S=K\cap\{n\cdot x=t\},\qquad
 S_0=K\cap\{n\cdot x=0\},\qquad n\ne0.
\]
For $p,q\in S$, convexity and central symmetry put $(p-q)/2$ in $K$,
and $n\cdot((p-q)/2)=0$.  Consequently
\[
 (S-S)/2\subseteq S_0.
\]
Identify the two parallel planes by translation.  The finite planar
construction gives a polygon $W\subseteq S_0$ having exactly the area
of $S$.  Hence
\[
 \boxed{\ \operatorname{area}(S)\le\operatorname{area}(S_0).\ }
\]
This argument includes supporting-plane sections, so no continuity at
a supporting face is being assumed.

\subsection{Arithmetic coordinates and the iPRA invariant}
For arithmetic inputs, an orthonormal frame is unnecessary.  Drop any
coordinate $j$ for which $n_j\ne0$.  The remaining two coordinates
provide an injective affine chart for both parallel planes; their
Euclidean areas have the same positive conversion factor
$\|n\|/|n_j|$.  Apply the preceding construction and area comparison to
the projected polygons.  All breakpoints, fiber endpoints, translated
coordinates, and trapezoid areas belong to the original ordered
arithmetic field.  The common positive conversion factor cancels.

The construction uses at most $2m$ distinct breakpoints and at most
$4m$ listed boundary vertices for $W$.  Sorting, merging, boundary-chain
traversal, and evaluating affine functions are primitive recursive.
A quantifier-free induction invariant for the summation loop is
\[
 2\sum_{i<k}U_i
   =\sum_{i<k}T_i+\sum_{i<k}T_{r-i-1}
 \qquad(0\le k\le r),                              \tag{CS5}
\]
where all sums are defined primitive-recursively on the finite lists.
The step is exactly (CS3); at $k=r$, reversing the finite sum gives
(CS4).  The convexity and containment checks are finite endpoint
inequalities and arithmetic midpoint identities.  This supplies a
direct iPRA proof of the central-section reduction for coded polygons,
without any excluded-middle decision about arbitrary real inputs or
any unbounded search.
